% Book 3, Year 1, chapter 5 -- dynamic memory.
% Provenance (never printed): sources/book3/BRIEFS.md, chapter 5.
% ledger: b1:malloc-spec, b1:lsan
\chapter{Dynamic Memory}\label{ch:b1:dynamic-memory}

A program must read the rainfall readings of a file, and the file may hold forty numbers or forty million. The \omterm{def:b1:pointers-and-arrays:array}{arrays}
of \cref{ch:b1:pointers-and-arrays} have a size fixed in the source; a local \omterm{def:b1:pointers-and-arrays:array}{array} lives in a frame of a stack
limited to 8~MiB; and both vanish when their function or the program ends. What the program needs is memory asked
for while it runs, in the amount the data require, kept for exactly as long as it is useful. C provides it, and
hands the programmer the duty to give it back. This chapter is about both halves of that bargain.

\begin{recall}
A running program's memory has regions: code, initialised and zero-initialised data, the \omterm{def:b1:c-functions-and-stack:call-stack}{call stack}, and the \omterm{def:b1:dynamic-memory:heap}{heap}
(\cref{ch:b1:compilation-pipeline}). A \omterm{def:b1:pointers-and-arrays:pointer}{pointer} holds an address and a type; \code{void *} holds an address of
anything; and \code{sizeof} gives sizes in bytes (\cref{ch:b1:pointers-and-arrays}).
\end{recall}

\section{\texttt{malloc} and \texttt{free}}\label{sec:b1:dynamic-memory:malloc}

\begin{definition}[Heap, dynamic allocation]\label{def:b1:dynamic-memory:heap}
The \emph{heap}\index{heap (memory)} is the region of a running program's memory from which blocks are handed out on
request while it runs. \emph{Dynamic allocation}\index{dynamic allocation} is the obtaining of such a block:
\code{malloc(n)} returns the address of a new block of \code{n} bytes, or the \omterm{def:b1:pointers-and-arrays:null}{null pointer} if it cannot. The block
exists until the program passes its address to \code{free}, whatever functions return in between.
\end{definition}

\omcode{code/bachelor-1/05-dynamic-memory/listings/alloc_basic.c}{4}{27}{An \omterm{def:b1:pointers-and-arrays:array}{array} whose length comes from the command
line}\label{lst:b1:dynamic-memory:basic}

\omoutput{code/bachelor-1/05-dynamic-memory/outputs/alloc_basic.txt}{Run with the argument 1000.}

Four habits are visible in these lines. The size is written \code{n * sizeof *squares}: the size of what the \omterm{def:b1:pointers-and-arrays:pointer}{pointer}
points to, which stays right if the type of \code{squares} changes. The result of \code{malloc}, a \code{void *}, is
stored without a \omterm{def:b1:c-types-and-control:cast}{cast}, since C converts it. The result is tested against \code{NULL} before use. And every
\code{malloc} has its \code{free}. \code{calloc(n, size)} allocates \code{n} elements and fills them with zero bytes;
\code{malloc} leaves the block's contents indeterminate.

\omcode{code/bachelor-1/05-dynamic-memory/listings/nullcheck.c}{2}{15}{Asking for too much}

\omoutput{code/bachelor-1/05-dynamic-memory/outputs/nullcheck.txt}{\code{malloc} cannot find nine exabytes and says
so.}\label{out:b1:dynamic-memory:null}

To change the size of a block, \code{realloc(p, n)} returns a block of \code{n} bytes holding the old contents
(as much as fits): the same block if it can grow in place, otherwise a new one, the old one being freed. If it fails it
returns \code{NULL} and leaves the old block untouched. So its result goes into a new variable: writing
\code{a = realloc(a, ...)} would lose the only address of the old block on failure
(\cref{exo:b1:dynamic-memory:5}).

\omcode{code/bachelor-1/05-dynamic-memory/listings/realloc_grow.c}{2}{27}{Growing a block
safely}\label{lst:b1:dynamic-memory:realloc}

\omoutput{code/bachelor-1/05-dynamic-memory/outputs/realloc_grow.txt}{The first four values survived the move.}

\begin{omfigure}
\begin{tikzpicture}[f/.style={draw, minimum width=30mm, minimum height=6mm, font=\small, anchor=north},
                    h/.style={draw, minimum height=6mm, font=\small, anchor=west}, >={Stealth[length=2mm]}]
\node[font=\small\bfseries] at (0,0.5) {call stack};
\node[f, fill=omCode!6] (s1) at (0,0) {\code{squares}};
\node[f, fill=omCode!6] (s2) at (0,-0.6) {\code{zeros}};
\node[f, fill=omCode!6] (s3) at (0,-1.2) {\code{n = 1000}};
\node[font=\small\bfseries] at (6.2,0.5) {heap};
\node[h, minimum width=40mm, fill=black!12, text=black!55] (b1) at (4.2,-0.3) {4000 bytes (freed)};
\node[h, minimum width=40mm, fill=omDef!10] (b2) at (4.2,-1.3) {4000 bytes, all zero};
\node[h, minimum width=16mm, fill=omDef!10] (b3) at (4.2,-2.3) {other block};
\draw[->, thick, black!45, dashed] (s1.east) -- (b1.west);
\draw[->, thick, omDef] (s2.east) -- (b2.west);
\node[font=\scriptsize, align=left, anchor=west] at (6.9,-2.3) {blocks of the C library\\and of other calls};
\end{tikzpicture}
\par\smallskip{\small Near the end of \cref{lst:b1:dynamic-memory:basic}. The \omterm{def:b1:pointers-and-arrays:pointer}{pointers} live in \code{main}'s frame;
the blocks they hold the addresses of live on the \omterm{def:b1:dynamic-memory:heap}{heap}. The first block has been freed: \code{squares} still holds
its address, but the block is no longer the program's.}
\end{omfigure}

\begin{omfigure}
\begin{tikzpicture}[>={Stealth[length=2mm]}, bar/.style={line width=5pt}]
\draw[->] (0,0) -- (11,0) node[right, font=\small] {time};
\foreach \x/\t in {0.5/\code{main}, 2.4/\code{f} called, 4.4/\code{malloc}, 6/\code{f} returns, 8.6/\code{free},
                   10.4/end}
  {\draw (\x,0.1) -- (\x,-0.1) node[below, font=\scriptsize] {\t};}
\draw[bar, omCode] (2.4,1.6) -- (6,1.6); \node[left, font=\small] at (0.3,1.6) {a local of \code{f}};
\draw[bar, omDef] (4.4,0.8) -- (8.6,0.8); \node[left, font=\small] at (0.3,0.8) {the heap block};
\node[font=\scriptsize, above, omDef] at (6.5,0.9) {outlives \code{f}};
\end{tikzpicture}
\par\smallskip{\small Two \omterm{def:b1:c-functions-and-stack:lifetime}{lifetimes}. A local of \code{f} ends when \code{f} returns; a block that \code{f} obtains
from \code{malloc} lasts until some part of the program frees it, possibly far from \code{f}. The programmer, not the
call structure, chooses when.}
\end{omfigure}

Compare with \cref{ch:b1:c-functions-and-stack}: a function could not return the address of a local, which died with
its frame; it can return the address of a block it allocated, which lives on. In exchange, the program must know
at every moment which blocks it still needs, and free the others exactly once.

\section{Errors and the tools that find them}\label{sec:b1:dynamic-memory:errors}

Manual memory management fails in four classic ways. None is detected by the language; each is detected, when a
run performs it, by the address sanitizer of \cref{ch:b1:pointers-and-arrays}.

\begin{definition}[Memory leak]\label{def:b1:dynamic-memory:leak}
A \emph{memory leak}\index{memory leak} is a \omterm{def:b1:dynamic-memory:heap}{heap} block that the program no longer frees because it no longer holds its
address anywhere, or no longer needs it but keeps it. A leaking program's memory grows with the time it runs.
\end{definition}

\omcode{code/bachelor-1/05-dynamic-memory/listings/leak.c}{2}{23}{A \omterm{def:b1:pointers-and-arrays:pointer}{pointer} overwritten before its block was
freed}\label{lst:b1:dynamic-memory:leak}

\omoutput[1-12]{code/bachelor-1/05-dynamic-memory/outputs/run_leak.txt}{At the end of the run, the leak detector lists
the blocks still allocated that nothing points to.}\label{out:b1:dynamic-memory:leak}

The address sanitizer includes a leak detector, LeakSanitizer, which runs when the program ends. It scans the memory
the program can still reach (globals, the stack, registers, the blocks reachable from them) for values that look like
addresses of \omterm{def:b1:dynamic-memory:heap}{heap} blocks; a block that nothing points to is reported. The two blocks of \code{item 1} and
\code{item 2} were lost when \code{s} was overwritten. The scan is conservative: a stale copy of an address left in a
dead stack slot makes its block look reachable, which is why this session builds at \code{-O1}, where no such copy
survives. The detector may miss a leak; it does not invent one.

\begin{definition}[Dangling pointer, use after free, double free]\label{def:b1:dynamic-memory:dangling}
After \code{free(p)}, every \omterm{def:b1:pointers-and-arrays:pointer}{pointer} that held the block's address is a \emph{dangling pointer}\index{dangling pointer}:
it still holds the number, but the block is no longer the program's. Reading or writing through it is a \emph{use
after free}\index{use after free}; passing it to \code{free} again is a \emph{double free}\index{double free}. Both are
\omterm{def:b1:c-types-and-control:ub}{undefined behaviour}.
\end{definition}

\omcode{code/bachelor-1/05-dynamic-memory/listings/uaf.c}{2}{14}{A read after
\code{free}}\label{lst:b1:dynamic-memory:uaf}

\omoutput[1-6,11-14,17-19]{code/bachelor-1/05-dynamic-memory/outputs/run_uaf.txt}{The report gives three stacks: the
bad read, the \code{free}, and the \code{malloc}.}\label{out:b1:dynamic-memory:uaf}

A use-after-free report tells the whole story of the block: where it was allocated (line 7), where it was freed
(line 11), and where it was used afterwards (line 12). Without the sanitizer, this program would most likely print 17,
since nobody has reused the block yet; the next \code{malloc} could reuse it, and then the read would return someone
else's data.

\omcode{code/bachelor-1/05-dynamic-memory/listings/double_free.c}{2}{11}{Two \omterm{def:b1:pointers-and-arrays:pointer}{pointers}, one block, two
\code{free}}\label{lst:b1:dynamic-memory:double}

\omoutput[1-6,12-14]{code/bachelor-1/05-dynamic-memory/outputs/run_double_free.txt}{The second \code{free} is
reported, with the first.}\label{out:b1:dynamic-memory:double}

\begin{omfigure}
\begin{tikzpicture}[c/.style={draw, minimum width=26mm, minimum height=6.5mm, font=\small},
                    p/.style={draw, minimum width=11mm, minimum height=6.5mm, font=\small\ttfamily},
                    >={Stealth[length=2mm]}]
\node[font=\small\bfseries] at (1.5,1.1) {before \code{free(a)}};
\node[p] (a1) at (0,0.4) {a}; \node[p] (b1) at (0,-0.5) {b};
\node[c, fill=omDef!10] (k1) at (3,0) {block of 8 \code{int}};
\draw[->, thick, omDef] (a1.east) -- (k1.west); \draw[->, thick, omDef] (b1.east) -- (k1.west);
\node[font=\small\bfseries] at (8.0,1.1) {after \code{free(a)}};
\node[p] (a2) at (6.5,0.4) {a}; \node[p] (b2) at (6.5,-0.5) {b};
\node[c, fill=black!12, text=black!55] (k2) at (9.5,0) {freed};
\draw[->, thick, omThm, dashed] (a2.east) -- (k2.west); \draw[->, thick, omThm, dashed] (b2.east) -- (k2.west);
\node[font=\scriptsize, omThm, align=center] at (8.0,-1.3) {both dangling: \code{*b} is a use after free,\\
\code{free(b)} a double free};
\end{tikzpicture}
\par\smallskip{\small Freeing a block through one \omterm{def:b1:pointers-and-arrays:pointer}{pointer} leaves every other copy of its address dangling. Nothing in
the \omterm{def:b1:pointers-and-arrays:pointer}{pointers} changes; only the block's status does, which is why the error is invisible in the source.}
\end{omfigure}

The fourth error, reading memory that was never given a value, is not an address error and the address sanitizer
does not see it. For a local variable the \omterm{def:b1:compilation-pipeline:compiler}{compiler} often does.

\omcode{code/bachelor-1/05-dynamic-memory/listings/uninit.c}{2}{12}{A total never set to
zero}\label{lst:b1:dynamic-memory:uninit}

\omoutput{code/bachelor-1/05-dynamic-memory/outputs/uninit.txt}{The warning of \code{-Wall} at
\code{-O2}.}\label{out:b1:dynamic-memory:uninit}

For a block from \code{malloc}, whose contents are indeterminate, no warning comes: the block may be filled by a
loop the \omterm{def:b1:compilation-pipeline:compiler}{compiler} cannot follow. Use \code{calloc}, or write every element before reading it. A debugger, the tool
that stops a running program and shows its variables, is the subject of a section of
\cref{ch:b1:c-modules-make-tests}.

\section{Who frees what}\label{sec:b1:dynamic-memory:convention}

Every block must be freed once, by someone, after its last use. In a program of a few hundred lines with blocks passed
between functions, this works only with a convention, written where every caller reads it: in the comment of the
interface. The C library itself follows one: a block from \code{malloc} is the caller's to free; a \omterm{def:b1:pointers-and-arrays:pointer}{pointer} returned by
\code{getenv} must never be freed.

\begin{method}[Writing down who frees]\label{met:b1:dynamic-memory:convention}
For each function that receives or returns a \omterm{def:b1:pointers-and-arrays:pointer}{pointer} to a \omterm{def:b1:dynamic-memory:heap}{heap} block, state in its interface comment:
\begin{enumerate}
\item whether a returned block belongs to the caller, who must free it (as for \code{malloc});
\item whether a \omterm{def:b1:pointers-and-arrays:pointer}{pointer} it receives is kept (the function takes charge of the block) or only borrowed during the call;
\item for a returned \omterm{def:b1:pointers-and-arrays:pointer}{pointer} into a block that the callee keeps (a borrowed \omterm{def:b1:pointers-and-arrays:pointer}{pointer}), until when it stays valid.
\end{enumerate}
Then give every structure a pair of functions, one that creates and one that destroys, and free in the destroying
function everything the structure holds.
\end{method}

The phone book of the weekend problem states its convention in its header; here it is, with the structure it
describes. A \code{struct} groups several variables under one name, accessed with \code{.} or, through a \omterm{def:b1:pointers-and-arrays:pointer}{pointer},
\code{->}; \cref{ch:b1:c-strings-structs-files} treats structures in full.

\omcode{code/bachelor-1/05-dynamic-memory/listings/phonebook.h}{2}{28}{An interface that says who
frees}\label{lst:b1:dynamic-memory:phonebook-h}

The \omterm{def:b1:pointers-and-arrays:pointer}{pointer} returned by \code{book_find} is borrowed, and the comment says until when: an addition may move the whole
\omterm{def:b1:pointers-and-arrays:array}{array} of contacts to a bigger block, and a removal frees a name. The Year~2 volume shows a language in which this
kind of convention is checked by the \omterm{def:b1:compilation-pipeline:compiler}{compiler}; in C it is checked by people and by sanitizers.

\section{A growable array written by hand}\label{sec:b1:dynamic-memory:growable}

To read an unknown number of values, keep them in a \omterm{def:b1:dynamic-memory:heap}{heap} block with room to spare, and move them to a bigger block when
it is full. The structure holds the block's address, the number of elements used, and the capacity of the block.

\omcode{code/bachelor-1/05-dynamic-memory/listings/intvec.c}{4}{34}{Pushing onto a growable \omterm{def:b1:pointers-and-arrays:array}{array} of
\code{int}}\label{lst:b1:dynamic-memory:intvec}

\omoutput{code/bachelor-1/05-dynamic-memory/outputs/read_all.txt}{Forty daily rainfalls of the town, in tenths of a
millimetre, read from standard input.}\label{out:b1:dynamic-memory:read-all}

\begin{definition}[Growth factor]\label{def:b1:dynamic-memory:growth}
A growable \omterm{def:b1:pointers-and-arrays:array}{array} whose capacity, when full, is multiplied by a fixed $r > 1$ has \emph{growth factor}\index{growth
factor (of an array)} $r$. The alternative is to add a fixed number $k$ of slots at each growth.
\end{definition}

Growth costs copies: moving from a full block of capacity $c$ copies $c$ elements. With growth by one slot, reaching
$n$ elements copies $1 + 2 + \dots + (n-1) = n(n-1)/2$ elements; with doubling, it copies $1 + 2 + 4 + \dots$ up to the
last power of two below $n$, fewer than $2n$. The program below counts both without allocating anything.

\omcode{code/bachelor-1/05-dynamic-memory/listings/copies.c}{2}{27}{Counting the copies of two growth
rules}\label{lst:b1:dynamic-memory:copies}

\omoutput{code/bachelor-1/05-dynamic-memory/outputs/copies.txt}{Half a trillion copies against a million, for a million
elements.}\label{out:b1:dynamic-memory:copies}

\begin{proposition}[Doubling copies fewer than $2n$ elements]\label{prop:b1:dynamic-memory:doubling}
A growable \omterm{def:b1:pointers-and-arrays:array}{array} that starts with capacity 1 and doubles its capacity when full copies fewer than $2n$ elements in all
while $n \geq 1$ elements are pushed.
\end{proposition}

\begin{proof}
The growths happen at capacities $1, 2, 4, \dots, 2^m$, where $2^m$ is the largest power of two below $n$ (a growth
happens when the \omterm{def:b1:pointers-and-arrays:array}{array} is full and one more element arrives, so its capacity is less than $n$). They copy
$1 + 2 + \dots + 2^m = 2^{m+1} - 1$ elements, and $2^m < n$ gives $2^{m+1} - 1 < 2n$. If $n = 1$, there is no growth
and no copy.
\end{proof}

So each push costs, on average over the whole sequence, fewer than two copies, however large $n$ is; doubling also
wastes at most half the block. \Cref{ch:b1:dynamic-arrays} turns this observation into a method, amortised analysis,
and calls the structure by its usual name.

\begin{omfigure}
\begin{tikzpicture}[c/.style={draw, minimum width=5.5mm, minimum height=5.5mm}]
\foreach \cap/\y in {1/0, 2/-0.8, 4/-1.6, 8/-2.4, 16/-3.2} {
  \node[font=\small, anchor=east] at (-0.3,\y) {capacity \cap};
  \pgfmathtruncatemacro{\old}{\cap/2}
  \foreach \k in {1,...,\cap} {
    \ifnum\k>\old \def\fillc{white}\else \def\fillc{omThm!25}\fi
    \node[c, fill=\fillc] at (\k*0.55,\y) {};
  }
}
\node[font=\scriptsize, align=left, anchor=west] at (9.4,-1.6) {shaded: elements\\copied from the\\previous block};
\end{tikzpicture}
\par\smallskip{\small Doubling from capacity 1 to 16. Each new block receives a copy of the full block before it, half
its size: $1 + 2 + 4 + 8 = 15$ copies to reach 16 elements, fewer than $2 \times 16$.}
\end{omfigure}

\section{Two-dimensional arrays of unknown size}\label{sec:b1:dynamic-memory:grids}

A grid of $h \times w$ numbers known only at run time can be stored in two ways. The first allocates a table of $h$
row \omterm{def:b1:pointers-and-arrays:pointer}{pointers}, then one block per row: \code{g[i][j]} works as for a static \omterm{def:b1:pointers-and-arrays:array}{array}, but there are $h + 1$ blocks to
allocate, check and free.

\omcode{code/bachelor-1/05-dynamic-memory/listings/grid_rows.c}{5}{30}{One block per
row}\label{lst:b1:dynamic-memory:rows}

\omoutput{code/bachelor-1/05-dynamic-memory/outputs/grid_rows.txt}{Four blocks for a $3 \times 5$ grid.}

\begin{definition}[Linearised array]\label{def:b1:dynamic-memory:linearised}
A \emph{linearised array}\index{linearised array} stores the elements of an $h \times w$ grid in a single block of
$h \cdot w$ elements, in \omterm{def:b1:pointers-and-arrays:row-major}{row-major order}: element $(i, j)$ is at index $i \cdot w + j$.
\end{definition}

\omcode{code/bachelor-1/05-dynamic-memory/listings/grid_flat.c}{2}{15}{The same grid in one
block}\label{lst:b1:dynamic-memory:flat}

\omoutput{code/bachelor-1/05-dynamic-memory/outputs/grid_flat.txt}{One allocation, one \code{free}.}

The linearised form makes one allocation and one \code{free}, keeps all elements contiguous (which the caches of the
processor reward, as the Year~2 volume measures), and lets the program pass the whole grid as one \omterm{def:b1:pointers-and-arrays:pointer}{pointer}. The table of
rows allows rows of different lengths, and swapping two rows by exchanging two \omterm{def:b1:pointers-and-arrays:pointer}{pointers}. Images, matrices and the
spreadsheet of \cref{ch:b1:pointers-and-arrays} are nearly always linearised.

\begin{safety}
Build and test every program that allocates with \code{-fsanitize=address}, which includes the leak detector. Free
every block exactly once, set a \omterm{def:b1:pointers-and-arrays:pointer}{pointer} to \code{NULL} after freeing it when it outlives the \code{free} (freeing
\code{NULL} does nothing), and never return or keep a \omterm{def:b1:pointers-and-arrays:pointer}{pointer} whose block someone else may free.
\end{safety}

\section{Exercises}\label{sec:b1:dynamic-memory:exercises}

\begin{exercise}[$\star$]\label{exo:b1:dynamic-memory:1}
Each function below contains one memory error. Name it, and say whether the address sanitizer would report it.
\omcode{code/bachelor-1/05-dynamic-memory/listings/exo_which.c}{5}{36}{Four errors}
\end{exercise}

\begin{exercise}[$\star$]\label{exo:b1:dynamic-memory:2}
How many bytes does \code{malloc(n * sizeof *r)} request for $n = 10$ in the program below? Predict, then explain the
size of \code{struct reading}, knowing that a \code{double} must start at a multiple of 8.
\omcode{code/bachelor-1/05-dynamic-memory/listings/exo_bytes.c}{2}{16}{The size of a request}
\end{exercise}

\begin{exercise}[$\star$]\label{exo:b1:dynamic-memory:3}
Predict the four numbers printed, and say which variable may no longer be used after the \code{realloc}.
\omcode{code/bachelor-1/05-dynamic-memory/listings/exo_realloc.c}{2}{20}{What survives a \code{realloc}}
\end{exercise}

\begin{exercise}[$\star$]\label{exo:b1:dynamic-memory:4}
One path through \code{sum_file} leaks. Which one? Fix it with one line.
\omcode{code/bachelor-1/05-dynamic-memory/listings/exo_where_leak.c}{5}{22}{Find the leak}
\end{exercise}

\begin{exercise}[$\star\star$]\label{exo:b1:dynamic-memory:5}
The function \code{grow} below is correct when memory is plentiful. What happens when \code{realloc} fails? Rewrite it.
\omcode{code/bachelor-1/05-dynamic-memory/listings/exo_realloc_bad.c}{2}{11}{The tempting \code{realloc}}
\end{exercise}

\begin{exercise}[$\star\star$]\label{exo:b1:dynamic-memory:6}
The loop at the end of the program frees a chain of three nodes. Read the report, explain the bug, and fix the loop.
\omcode{code/bachelor-1/05-dynamic-memory/listings/exo_chain.c}{5}{25}{Freeing a chain}
\omoutput[1-6]{code/bachelor-1/05-dynamic-memory/outputs/run_chain.txt}{The first lines of the report.}
\end{exercise}

\begin{exercise}[$\star\star$]\label{exo:b1:dynamic-memory:7}
Write \code{char *my_strdup(const char *s)}, which returns a new \omterm{def:b1:dynamic-memory:heap}{heap} block holding a copy of the string \code{s}
(its characters and the final zero byte), or \code{NULL}. State its convention for freeing.
\end{exercise}

\begin{exercise}[$\star\star$]\label{exo:b1:dynamic-memory:8}
Write \code{int **grid2d_new(size_t h, size_t w)}, which allocates a zero-filled grid as a table of rows, and
\code{grid2d_free}. If an allocation fails halfway, \code{grid2d_new} must return \code{NULL} without leaking.
\end{exercise}

\begin{exercise}[$\star\star$]\label{exo:b1:dynamic-memory:9}
A growable \omterm{def:b1:pointers-and-arrays:array}{array} starts with capacity $k$ and adds $k$ slots when full. Show that pushing $n$ elements copies
$k\,m(m+1)/2$ elements, where $m = \lfloor (n-1)/k \rfloor$, roughly $n^2/(2k)$. Check it against the \code{+16} row of
\cref{out:b1:dynamic-memory:policies} for $n = 1000$.
\end{exercise}

\begin{exercise}[$\star\star\star$]\label{exo:b1:dynamic-memory:10}
Prove \cref{prop:b1:dynamic-memory:doubling} again for an \omterm{def:b1:pointers-and-arrays:array}{array} that starts at capacity $c_0 \geq 1$, and show that the
bound $2n$ cannot be replaced by $(2 - \varepsilon)n$ for any $\varepsilon > 0$.
\end{exercise}

\begin{exercise}[$\star\star\star$]\label{exo:b1:dynamic-memory:11}
A program allocates and frees many blocks of one size. Write a pool: one \code{malloc} of \code{count} blocks at the
start, then \code{pool_get} and \code{pool_put} in constant time, keeping the free blocks in a chain that runs through
their own first bytes. Why must each block be at least as large as a \omterm{def:b1:pointers-and-arrays:pointer}{pointer}?
\end{exercise}

\begin{exercise}[$\star\star\star$]\label{exo:b1:dynamic-memory:12}
During a \code{realloc} that moves, the old and the new block exist at once. For an \omterm{def:b1:pointers-and-arrays:array}{array} that doubles from capacity 1,
prove that the number of slots allocated at the same moment never exceeds $3n$, where $n$ is the number of elements
after the push that triggered the growth.
\end{exercise}

\section{Problem: The leaky phone book}\label{sec:b1:dynamic-memory:problem}

\begin{problem}[{Weekend problem --- The leaky phone book}]\label{pb:b1:dynamic-memory:1}
The phone book of \cref{lst:b1:dynamic-memory:phonebook-h} stores contacts in a growable \omterm{def:b1:pointers-and-arrays:array}{array}, under the convention
written in its header. Its first version, \code{phonebook_v0.c}, and a short session, \code{pbdemo.c}, contain five
memory errors, numbered in comments. The session was built with the address sanitizer in a mode that reports an error
and goes on, so that one run shows them all; the report was saved, and its summary lines extracted. A second program,
\code{policies.c}, adds contacts to the corrected book under three growth policies.
\omcodes{code/bachelor-1/05-dynamic-memory/listings/phonebook_v0.c}{22-44,54-70}{The first version: adding, removing,
destroying}\label{lst:b1:dynamic-memory:v0}
\omcode{code/bachelor-1/05-dynamic-memory/listings/pbdemo.c}{6}{22}{The session}\label{lst:b1:dynamic-memory:pbdemo}
\omoutput{code/bachelor-1/05-dynamic-memory/outputs/problem.txt}{The start of the first report, then every summary
line and the session's own output.}\label{out:b1:dynamic-memory:problem}
\omoutput{code/bachelor-1/05-dynamic-memory/outputs/policies.txt}{Copies, final capacity and unused slots for three
policies.}\label{out:b1:dynamic-memory:policies}

\medskip\noindent\textbf{Part I --- Reading the reports.}
\begin{enumerate}
\item The first report is a write of 16 bytes into a block of 15. Which call writes, which line of \code{book_add}
made the block, and why is it one byte short?
\item Why does the summary line of that report name \code{strcpy} rather than \code{book_add}, and how does the start of
the report recover the line of \code{book_add}?
\item The second summary is a read by \code{strcmp}. Which function calls it, and why is it the same error as the first?
\item The third summary is a \omterm{def:b1:dynamic-memory:dangling}{use after free} at line 12 of \code{pbdemo.c}. Which block was freed, by which call, and
why?
\item The fourth summary is in \code{book_remove}. With 4 contacts in a block of capacity 4, which element does
\code{b->items[b->len]} designate?
\item The fifth summary is at line 17 of \code{pbdemo.c}. What does \code{who} point to, and who freed it?
\item The leak report counts 41 bytes in 3 blocks. Identify the three blocks, and check the total with the lengths of
the names and error 1.
\end{enumerate}

\medskip\noindent\textbf{Part II --- Fixing under the convention.}
\begin{enumerate}[resume]
\item Fix error 1. Why is \code{strlen(name) + 1} the right size?
\item Fix error 2, and say which contact the erroneous line lost.
\item Fix error 3 so that \code{book_destroy} frees everything the book holds, in the right order.
\item Error 4 is in the caller. Which sentence of the convention did it break? Rewrite the session so that it is
correct.
\item Same question for error 5.
\item With all five fixed, which blocks exist just before \code{book_destroy}, and how many \code{free} calls does it
make?
\item Why does \code{book_add} copy the name instead of keeping the caller's \omterm{def:b1:pointers-and-arrays:pointer}{pointer}? Give a caller for whom keeping
it would be a bug.
\end{enumerate}

\medskip\noindent\textbf{Part III --- Growth policies.}
\begin{enumerate}[resume]
\item For the policy \code{+16} and $n = 1000$, recompute the 31\,248 copies of \cref{out:b1:dynamic-memory:policies}.
\item For doubling and $n = 1000$: list the capacities reached and recompute 1023 copies and 24 unused slots.
\item The policy x1.5 computes the new capacity as $\lceil 1.5c \rceil$. Give its first seven capacities, from 1.
\item For $n = 100\,000$, by what factor does \code{+16} copy more than doubling? Explain the order of magnitude.
\item Why does the policy \code{+16} have 0 unused slots for $n = 10\,000$ and $n = 100\,000$?
\end{enumerate}

\medskip\noindent\textbf{Part IV --- The memory budget.}
\begin{enumerate}[resume]
\item Let a policy replace a full capacity $c$ by a capacity $C$ with $rc \leq C < rc + 1$, for a fixed $r > 1$ (doubling,
and x1.5 rounded up, are such policies). Prove that pushing $n$ elements copies fewer than $rn/(r-1)$ elements.
\item Prove that the unused slots after the last growth are fewer than $(r-1)n$.
\item Check both bounds on the rows x2 and x1.5 of \cref{out:b1:dynamic-memory:policies}.
\item A contact takes 16 bytes. For $n = 100\,000$, how many bytes does each policy leave unused? Which policy would you
choose for a phone with little memory, and which for a server that adds contacts all day?
\item State the named result of the problem, \emph{the growth-factor trade-off}.
\end{enumerate}
\end{problem}
